\documentclass{article}
\usepackage[UTF8]{ctex}
\usepackage{amsmath,amssymb,amsfonts}
\usepackage{enumitem}
\usepackage{geometry}
\geometry{a4paper, margin=1in}
% === 错题本排版优化（不参与题目标识） ===
\setlength{\parindent}{0pt}
\setlength{\parskip}{0.35em plus 0.10em minus 0.08em}
\setlength{\emergencystretch}{3em}
\setlist{itemsep=0.25em, topsep=0.40em, parsep=0pt}
\allowdisplaybreaks[2]
% === 错题本排版优化结束 ===
\begin{document}
\section*{华中师范大学2024年高等代数真题}
\section*{129．华中师范大学2024年高等代数试卷}
一．填空题。\\

\subsection{题目1}
1．设 \(A, B\) 均为 \(n\) 阶方阵，\(r(A)=1, A B=O\) ，则 \(B\) 的秩最大值是 \(\_\_\_\_\)．

\subsection{题目2}
2．设 \(A=\left(\begin{array}{rrr}1 & * & * \\ & 4 & * \\ & & 6\end{array}\right)\) ，则 \(A^{2}\) 的迹是 \(\_\_\_\_\)

\subsection{题目3}
3．设 \(A=\left(\begin{array}{cc}0 & -1 \\ 2 & 3\end{array}\right)\) ，则 \(A^{2023}=\) \(\_\_\_\_\)．\\
\(\_\_\_\_\)

\subsection{题目4}
4．若 \(V=\left\{X=\left(x_{1}, x_{2}, x_{3}\right)^{\mathrm{T}} \mid x_{1}+x_{2}+x_{3}=0\right\}\) ，则 \(\operatorname{dim} V=\) \(\_\_\_\_\)．

\subsection{题目5}
5．矩阵 \(\left(\begin{array}{lll}1 & 1 & 1 \\ 1 & 1 & 1 \\ 1 & 1 & 1\end{array}\right)\) 的若尔当标准形为 \(\_\_\_\_\)．\\
\(\_\_\_\_\)

\subsection{题目6}
6．设 \(f(\lambda), g(\lambda)\) 为数域 \(\mathbb{C}\) 上的多项式，且 \((f(\lambda), g(\lambda))=1\) ，则 \(\lambda\) —矩阵 \(\left(\begin{array}{ll}f(\lambda) & \\ & g(\lambda)\end{array}\right)\) 的等价标准形为 \(\_\_\_\_\)．

二．解答题。\\

\subsection{题目7}
7．计算行列式 \(D=\left|\begin{array}{cccc}a & 0 & b & 0 \\ 0 & c & 0 & d \\ -b & 0 & -a & 0 \\ 0 & -d & 0 & -c\end{array}\right|\) 的值。

\subsection{题目8}
8．解答如下问题：

\begin{itemize}
  \item[（1）] 证明：\(x^{3}-2\) 作为有理系数多项式不可约．
  \item[（2）] 求有理数 \(a, b, c\) ，使得 \(a \sqrt[3]{4}+b \sqrt[3]{2}+c=\frac{1}{\sqrt[3]{4}+\sqrt[3]{2}+4}\) 成立。
\end{itemize}

\subsection{题目9}
9．设 \(A\) 为 \(n\) 阶实矩阵，证明：\(A\) 的前 \(r\) 个列向量是 \(A\) 的列向量组的极大线性无关组的充分必要条件是 \(A^{\mathrm{T}} A\) 的前 \(r\) 列是 \(A^{\mathrm{T}} A\) 的列向量组的一个极大线性无关组。

\subsection{题目10}
10．二次型 \(q\left(x_{1}, x_{2}, x_{3}\right)=x_{1}^{2}+2 x_{2}^{2}+x_{3}^{2}+2 a x_{1} x_{2}+2(1+a) x_{1} x_{3}\) 为正定二次型，求 \(a\) 的取值范围．

\subsection{题目11}
11．设 \(V\) 为数域 \(\mathbb{C}\) 上的 \(n\) 维线性空间，\(\varphi\) 为 \(V\) 上的不可逆的线性变换．证明：存在 \(V\) 上非零线性变换 \(\psi\) ，使得 \(\psi \varphi=\varphi \psi=0\) 。

\subsection{题目12}
12．设 \(Q\) 为 \(n\) 阶复矩阵，且满足 \(Q \bar{Q}^{\mathrm{T}}=E_{n \times n}\) ．

\begin{itemize}
  \item[（1）] 证明：\(Q\) 的所有特征值的模长为 1 。
  \item[（2）] 举例说明是否存在 2 阶复矩阵 \(P\) ，其特征值模长均为 1 ，但是 \(P \bar{P}^{\mathrm{T}} \neq E_{2 \times 2}\) ．
\end{itemize}
\end{document}
